%\titlespacing{\section}{0pt}{12pt plus 4pt minus 2pt}{-6pt plus 2pt minus 2pt}

% Kopf- und Fusszeile
\renewcommand{\sectionmark}[1]{\markright{#1}}
\renewcommand{\leftmark}{\rightmark}
\pagestyle{fancy}
\lhead{}
\chead{}
\rhead{\thesection\space\contentsname}
\lfoot{Implementierung von Brettspielen am Beispiel ReversiXT -- SS \the\year}
\cfoot{}
\rfoot{\ \linebreak Seite \thepage}
\renewcommand{\headrulewidth}{0.4pt}
\renewcommand{\footrulewidth}{0.4pt}

% Vorspann
\renewcommand{\thesection}{\Roman{section}}
\renewcommand{\theHsection}{\Roman{section}}
\pagenumbering{Roman}

\newcommand{\folgen}[1]{
\ensuremath
#1
}

\newcommand{\MyTitlepage}[5][\empty]{
\thispagestyle{empty}
\begin{center}
	\includegraphics[scale=0.2]{pics/oth-logo}\\
	\vspace*{2cm}
	\Large
	\textbf{Fakultät}\\
	\textbf{Informatik und Mathematik}\\
	\vspace*{2cm}
	\Huge
	\textbf{Projektbericht}\\
	\vspace*{0.5cm}
	\large
	zum Wahlpflichtfach im SS \the\year\\
	\vspace*{1cm}
	\textbf{Implementierung von Brettspielen am Beispiel ReversiXT}\\
	\vspace*{1cm}
	\includegraphics[height=6cm]{#1}
	\vfill
	\normalsize
	%\newcolumntype{x}[1]{>{\raggedleft\arraybackslash\hspace{0pt}}p{#1}}
	\begin{tabular}{rl}%{6cm}p{7.5cm}}
	    \rule{0mm}{5ex}\textbf{Gruppe:} & #2 \\
		\rule{0mm}{5ex}\textbf{Autoren:} & \hspace*{-0.5em}\begin{tabular}[t]{r}#3\end{tabular} \\ 
		\rule{0mm}{5ex}\textbf{Leiter:} & Prof.\ Dr.\ rer.\ nat.\ Carsten Kern \\
		\rule{0mm}{5ex}\textbf{Abgabedatum:} & #4 \\ 
	\end{tabular} 
\end{center}
\pagebreak
}

% Entspricht Symbol
\newcommand{\corresponds}{\ensuremath{\mathrel{\widehat{=}}} }